% META: {"id":"1SPE-DERGLOBAL-ME-004","chapitre":"1SPE-DERIVATION-GLOBAL","type_objet":"methode","capacites_codes":["C4"],"status":"generated"}
\begin{fichemethode}{M4}{Déterminer les extremums d'une fonction}
\quandUtiliser{J'utilise cette méthode lorsque l'énoncé demande de trouver un maximum ou un minimum local (ou global sur un intervalle) d'une fonction dérivable.}
\pasApas{\begin{enumerate}
  \item Je précise le domaine de travail, ses bornes incluses ou exclues,
        et les intervalles de dérivabilité.
  \item Je calcule $f'(x)$.
  \item Je résous $f'(x)=0$ pour trouver les candidats intérieurs.
  \item J'étudie le signe de $f'$ et les variations, y compris les éventuelles
        portions constantes, pour déterminer les extremums locaux.
  \item Pour un extremum global, je compare les valeurs aux candidats et aux
        bornes incluses en utilisant l'ensemble du tableau de variations.
\end{enumerate}}
\exempleRedige{Déterminons les extremums de $f(x)=-x^3+3x^2-1$ sur $\mathbb{R}$. \commentaireMarge{Dériver.} $f'(x)=-3x^2+6x=-3x(x-2)$. \commentaireMarge{Trouver les annulateurs.} $f'(x)=0 \Leftrightarrow x=0$ ou $x=2$. \commentaireMarge{Étudier le signe de $f'$.} $f'(x)<0$ sur $\left]-\infty;0\right[$, $f'(x)>0$ sur $\left]0;2\right[$, $f'(x)<0$ sur $\left]2;+\infty\right[$. \commentaireMarge{Conclure sur les extremums.} $f'$ change de $-$ à $+$ en $0$ : minimum local $f(0)=-1$. $f'$ change de $+$ à $-$ en $2$ : maximum local $f(2)=-8+12-1=3$.
Ces extremums ne sont pas globaux : $f(-2)=19>3$ et $f(4)=-17<-1$.
Comme tout extremum global sur $\mathbb R$ serait aussi local, il n'y en a aucun.}
\pieges{\begin{itemize}
  \item Conclure à un extremum à partir de la seule égalité $f'(x)=0$.
  \item Oublier une borne incluse lors de la recherche d'un extremum global.
  \item Confondre maximum et minimum en lisant mal le tableau de variations.
  \item Oublier de calculer la valeur $f(x_0)$ qui définit l'extremum.
\end{itemize}}
\verifier{L'annulation de la dérivée ne suffit pas : $x^3$ n'a pas d'extremum
en $0$. À l'inverse, une fonction constante a des extremums non stricts sans
changement de signe. Je justifie la conclusion à partir des variations.}
\sEntrainer{\refExos{M4}}
\end{fichemethode}
